%------DEPENDENCIES------
% \usepackage{tikz}
% \usepackage{graphicx}
% \usepackage{xparse}
%------------------------

%----SHOWCOORDS

\newcommand{\showcoords}{
\draw[step=0.1,gray!30,very thin] (0,0) grid (1,1);
\draw[step=0.5,gray!70,thin] (0,0) grid (1,1);

\foreach \x in {0,0.1,...,1}
  \node[below,scale=0.5] at (\x,0) {\x};

\foreach \y in {0,0.1,...,1}
  \node[left,scale=0.5] at (0,\y) {\y};
}

% \newcommand{\showcoords}{
% \draw[step=0.1,gray!40,very thin] (0,0) grid (1,1);

% \foreach \x in {0,0.1,...,1}
%   \draw[gray] (\x,0) -- (\x,-0.02) node[below,scale=0.5] {\x};

% \foreach \y in {0,0.1,...,1}
%   \draw[gray] (0,\y) -- (-0.02,\y) node[left,scale=0.5] {\y};
% }

%---\coordprobe
\newcommand{\coordprobe}{
\draw[step=0.1,gray!30,very thin] (0,0) grid (1,1);

\foreach \x in {0,0.1,...,1}
  \node[below,scale=0.5] at (\x,0) {\x};

\foreach \y in {0,0.1,...,1}
  \node[left,scale=0.5] at (0,\y) {\y};

\path
(0,0)
coordinate[pos=0] (dummy);

\foreach \p in {(0,0)}{
\node[opacity=0] at \p {};
}
}

%----\crosshair

\newcommand{\crosshair}[2]{
\draw[red] (#1-0.01,#2) -- (#1+0.01,#2);
\draw[red] (#1,#2-0.01) -- (#1,#2+0.01);
\node[above right,scale=0.6] at (#1,#2) {(#1,#2)};
}

%-----\imageannot

\newcommand{\imageannot}[2]{
\begin{tikzpicture}

\node[anchor=north west,inner sep=0] (img)
{\includegraphics[width=\textwidth]{#1}};

\begin{scope}
[x={(img.north east)},y={(img.south west)}]

#2

\end{scope}

\end{tikzpicture}
}

\newcommand{\imageannotdebug}[2]{
\begin{tikzpicture}

\node[anchor=north west,inner sep=0] (img)
{\includegraphics[width=\textwidth]{#1}};

\begin{scope}
[x={(img.north east)},y={(img.south west)}]

\showcoords

#2

\end{scope}

\end{tikzpicture}
}


\newcommand{\highlightfrag}[5][yellow]{
\fill[#1,opacity=0.35]
(#2,#3) rectangle (#4,#5);

\draw[#1,line width=0pt]
(#2,#3) rectangle (#4,#5);
}

\let\highlightmove\highlightfrag

% \newcommand{\imageannot}[3][false]{
% \begin{tikzpicture}

% \node[anchor=north west,inner sep=0] (img)
% {\includegraphics[width=\textwidth]{#2}};

% \begin{scope}
% [x={(img.north east)},y={(img.south west)}]

% \ifx#1true
% \showcoords
% \fi

% #3

% \end{scope}

% \end{tikzpicture}
% }

% \newcommand{\imageannot}[2]{
% \begin{tikzpicture}
% \node[anchor=north west,inner sep=0] (img)
% {\includegraphics[width=\textwidth]{#1}};

% \begin{scope}
% [x={(img.north east)},y={(img.south west)}]

% #2

% \end{scope}
% \end{tikzpicture}
% }

% \newcommand{\imageannot}[2]{
% \begin{tikzpicture}
% \node[anchor=south west,inner sep=0] (img) at (0,0)
% {\includegraphics[width=\textwidth]{#1}};

% \begin{scope}[x={(img.south east)},y={(img.north west)}]
% #2
% \end{scope}

% \end{tikzpicture}
% }

%---\circlefrag{}{}{}

% \newcommand{\circlefrag}[3]{
% \draw[red, line width=2pt] (#1,#2) circle (#3);
% }

%------ellipsefrag

\NewDocumentCommand{\ellipsefrag}{O{red} O{0} m m m m}{
\draw[#1,line width=1.5pt,
      rotate around={#2:({(#3+#5)/2},{(#4+#6)/2})}]
({(#3+#5)/2},{(#4+#6)/2})
ellipse
({(#5-#3)/2} and {(#6-#4)/2});
}

% \newcommand{\ellipsefrag}[5][red]{
% \draw[#1,line width=2pt]
% ({(#2+#4)/2},{(#3+#5)/2})
% ellipse
% ({(#4-#2)/2} and {(#5-#3)/2});
% }

% \newcommand{\ellipsefrag}[4]{
% \draw[red,line width=2pt]
% ({(#1+#3)/2},{(#2+#4)/2})
% ellipse
% ({(#3-#1)/2} and {(#4-#2)/2});
% }

%------ovalfrag

\NewDocumentCommand{\ovalfrag}{O{blue} O{0} O{0.2} m m m m}{
\fill[draw=none,#1,fill opacity=#3,
      rotate around={#2:({(#4+#6)/2},{(#5+#7)/2})}]
({(#4+#6)/2},{(#5+#7)/2})
ellipse
({(#6-#4)/2} and {(#7-#5)/2});
}

% \NewDocumentCommand{\ovalfrag}{O{blue} O{0} O{0.2} m m m m}{
% \fill[#1,fill opacity=#3,
%       rotate around={#2:({(#4+#6)/2},{(#5+#7)/2})}]
% ({(#4+#6)/2},{(#5+#7)/2})
% ellipse
% ({(#6-#4)/2} and {(#7-#5)/2});
% }

% \NewDocumentCommand{\ovalfrag}{O{red!20} O{0} m m m m}{
% \fill[#1,
%       rotate around={#2:({(#3+#5)/2},{(#4+#6)/2})}]
% ({(#3+#5)/2},{(#4+#6)/2})
% ellipse
% ({(#5-#3)/2} and {(#6-#4)/2});
% }

% \NewDocumentCommand{\ovalfrag}{O{red} O{0} m m m m}{
% \fill[#1,
%       rotate around={#2:({(#3+#5)/2},{(#4+#6)/2})}]
% ({(#3+#5)/2},{(#4+#6)/2})
% ellipse
% ({(#5-#3)/2} and {(#6-#4)/2});
% }

%----rectfrag{}{}{}{}

\NewDocumentCommand{\rectfrag}{O{red} O{0} m m m m}{
\draw[#1,line width=1.0pt,rounded corners=3pt,
      rotate around={#2:({(#3+#5)/2},{(#4+#6)/2})}]
(#3,#4) rectangle (#5,#6);
}

% \newcommand{\rectfrag}[5][red]{
% \draw[#1,line width=1.5pt,rounded corners=3pt]
% (#2,#3) rectangle (#4,#5);
% }

% \newcommand{\rectfrag}[5][red]{
% \draw[#1,line width=2pt,rounded corners=3pt]
% (#2,#3) rectangle (#4,#5);
% }

% \newcommand{\rectfrag}[4]{
% \draw[red,line width=2pt,rounded corners=3pt]
% (#1,#2) rectangle (#3,#4);
% }

% \newcommand{\rectfrag}[4]{
% \draw[red,line width=2pt] (#1,#2) rectangle (#3,#4);
% }

%---\arrowfrag{}{}{}{}

% \newcommand{\arrowfrag}[4]{
% \draw[red,line width=2pt,->]
% (#1,#2) -- (#3,#4);
% }

\NewDocumentCommand{\arrowfrag}{O{red} m m m m}{
\draw[#1,line width=2pt,-{Latex[length=4mm]}]
(#2,#3) -- (#4,#5);
}

% \newcommand{\arrowfrag}[4]{
% \draw[red,line width=2pt,-{Latex[length=4mm]}]
% (#1,#2) -- (#3,#4);
% }

%---\textfrag

\newcommand{\textfrag}[5][black]{
\node[#1,anchor=center,scale=#4,inner sep=0pt] at (#2,#3) {#5};
}

% \newcommand{\textfrag}[5][black]{
% \node[#1,anchor=center,scale=#4,inner sep=0pt] at (#2,#3) {#5};
% }

% \newcommand{\textfrag}[5][black]{
% \node[#1,anchor=north west,scale=#4] at (#2,#3) {#5};
% }

% \newcommand{\textfrag}[4]{
% \node[red,anchor=north west,scale=#3] at (#1,#2) {#4};
% }

%---\markerfrag



\newcommand{\markerfrag}[4][red]{
\node[
circle,
draw=#1,
fill=white,
line width=1.5pt,
inner sep=1.5pt,
font=\bfseries\small,
text=#1
] at (#2,#3) {#4};
}

\newcommand{\markerfragdark}[4][red]{
\node[
circle,
fill=#1,
text=white,
font=\bfseries\small,
inner sep=1.8pt
] at (#2,#3) {#4};
}

%---------
% \usetikzlibrary{calc}

\newcommand{\rulefrag}[5][black]{%
  % main line
  \draw[#1] (#2,#3) -- (#4,#5);
  % first end cap
  \draw[#1]
    ($ (#2,#3)!0.08cm!90:(#4,#5) $) --
    ($ (#2,#3)!0.08cm!-90:(#4,#5) $);
  % second end cap
  \draw[#1]
    ($ (#4,#5)!0.08cm!90:(#2,#3) $) --
    ($ (#4,#5)!0.08cm!-90:(#2,#3) $);
}

%-----

\newlength{\pinwidth}

\newcommand{\pinfrag}[5][0]{%
\pgfmathsetlength{\pinwidth}{#4*0.9cm}%
\node[anchor=center,inner sep=0pt] at (#2,#3)
{\includegraphics[width=\pinwidth,angle=#1]{#5}};
}

% \newcommand{\pinfrag}[5][0]{%
% \node[anchor=center,inner sep=0pt] at (#2,#3)
% {\includegraphics[width=\dimexpr#4*0.9cm\relax,angle=#1]{#5}};
% }

% \newcommand{\pinfrag}[5][0]{%
% \node[anchor=center] at (#2,#3)
% {\includegraphics[width=#4*0.9cm,angle=#1]{#5}};
% }

% \newcommand{\pinfrag}[5][0]{
% \node[
%     anchor=center,
%     rotate=#1,
%     scale=#4
% ] at (#2,#3)
% {\includegraphics{#5}};
% }

% \newcommand{\pinfrag}[5][0]{
% \node[anchor=center] at (#2,#3)
% {\includegraphics[scale=#4,angle=#1]{#5}};
% }


%--usage---

% \rulefrag[red]{0.2}{0.3}{0.8}{0.3}

%---- USAGE ----
% \imageannot{analysis.png}{
%   \circlefrag{0.25}{0.70}{0.05}
%   \circlefrag{0.55}{0.35}{0.08}
% }

% DEBUG MODE
% \imageannot[true]{analysis.png}{
% }

%===========================================================
% EXAMPLE USAGE / CHEAT SHEET
% (All examples commented out – purely for documentation)
%===========================================================

%-----------------------------------------------------------
% BASIC IMAGE ANNOTATION
%-----------------------------------------------------------
% Places an image and allows drawing annotations in normalized
% coordinates (0–1 horizontally and vertically).
%
% \imageannot{path/to/image.png}{
%   % annotation commands go here
% }

% Example:
% \imageannot{analysis.png}{
%   \rectfrag{0.2}{0.3}{0.5}{0.6}
% }

%-----------------------------------------------------------
% DEBUG MODE WITH COORDINATE GRID
%-----------------------------------------------------------
% Shows a grid with coordinates so you can measure where
% to place annotations.

% \imageannotdebug{analysis.png}{
% }

% You can also manually show the coordinate grid:
%
% \imageannot{analysis.png}{
%   \showcoords
% }

%-----------------------------------------------------------
% HIGHLIGHT RECTANGLE
%-----------------------------------------------------------
% Semi-transparent highlight box.
%
% Syntax:
% \highlightfrag[color]{x1}{y1}{x2}{y2}
%
% Default color: yellow

% Example:
% \imageannot{analysis.png}{
%   \highlightfrag{0.2}{0.3}{0.6}{0.5}
% }

% Custom color:
% \imageannot{analysis.png}{
%   \highlightfrag[green]{0.2}{0.3}{0.6}{0.5}
% }

% highlightmove is currently an alias:
%
% \highlightmove{0.2}{0.3}{0.6}{0.5}

%-----------------------------------------------------------
% RECTANGLE FRAME
%-----------------------------------------------------------
% Draws a red rounded rectangle.

% \imageannot{analysis.png}{
%   \rectfrag{0.2}{0.3}{0.6}{0.5}
% }

%-----------------------------------------------------------
% ELLIPSE / CIRCLE HIGHLIGHT
%-----------------------------------------------------------
% Draws an ellipse that fits inside the rectangle defined
% by two points.

% \imageannot{analysis.png}{
%   \ellipsefrag{0.3}{0.4}{0.6}{0.7}
% }

%-----------------------------------------------------------
% ARROW
%-----------------------------------------------------------
% Draws an arrow from one coordinate to another.

% \imageannot{analysis.png}{
%   \arrowfrag{0.2}{0.2}{0.5}{0.5}
% }

%-----------------------------------------------------------
% TEXT LABEL
%-----------------------------------------------------------
% Places text on the image.
%
% Syntax:
% \textfrag{x}{y}{scale}{text}

% \imageannot{analysis.png}{
%   \textfrag{0.3}{0.8}{1}{Important motif}
% }

%-----------------------------------------------------------
% MARKERS (NUMBERED CIRCLES)
%-----------------------------------------------------------
% White circle with red outline.

% \imageannot{analysis.png}{
%   \markerfrag{0.2}{0.7}{1}
%   \markerfrag{0.5}{0.4}{2}
% }

% Filled red circle with white text.

% \imageannot{analysis.png}{
%   \markerfragdark{0.2}{0.7}{A}
% }

%-----------------------------------------------------------
% CROSSHAIR FOR PRECISE COORDINATES
%-----------------------------------------------------------
% Draws a small crosshair and prints the coordinate.

% \imageannot{analysis.png}{
%   \crosshair{0.5}{0.5}
% }

%-----------------------------------------------------------
% MULTIPLE ANNOTATIONS TOGETHER
%-----------------------------------------------------------
% \imageannot{analysis.png}{
%   \highlightfrag{0.2}{0.3}{0.6}{0.5}
%   \arrowfrag{0.7}{0.7}{0.5}{0.5}
%   \textfrag{0.65}{0.75}{0.8}{Cadence}
%   \markerfrag{0.2}{0.3}{1}
% }

%-----------------------------------------------------------
% WORKFLOW TIP
%-----------------------------------------------------------
% 1. Use \imageannotdebug to see coordinates.
% 2. Note positions you want.
% 3. Replace with \imageannot and annotations.
%-----------------------------------------------------------

\newcommand{\parafrag}[5][\normalsize]{%
  \node[
    anchor=north west,
    text width=#4\textwidth,
    align=left
  ] at (#2,#3) {{#1 #5}};
}

\tikzset{
  whiteout/.style={fill=white, draw=none, opacity=1}
}

\renewcommand{\highlightfrag}[5][yellow]{%
  \path[fill, draw, line width=0pt, fill opacity=0.35, #1]
    (#2,#3) rectangle (#4,#5);
}

\let\highlightmove\highlightfrag